% Chapter 5. Every number is a macro (generated/numbers.tex) or a \needsclaim box.
\chapter{The Value of Thinking While Listening}
\label{ch:value}

Chapter~\ref{ch:cost} priced the thinker's speculation. This chapter asks what
it buys. The answer is a negative, and the aim is to state it precisely: on
this task a reply draft is rarely usable at any point where it can still
finish, and when it becomes usable, little time is left. The two triggers we
evaluated could not be shown to tell the cases apart; whether any trigger
could remains open.

\section{What a usable draft is}
\label{sec:usable}

A draft saves time only if what it would say first is what the final reply
says first. We call a draft usable if its first synthesized chunk is a word
prefix of the first chunk of the reply generated from the full transcript.
Matches that contain no digit and share no content word with the problem are
reported separately. The break-even usability follows from the measured costs: a draft
is worth issuing only if the probability that it is usable exceeds
$p^\ast = c / (c + s)$. Here $c$ is the overlap cost of a draft still running
at the endpoint, and $s$ the time a usable draft saves. On this pipeline
$p^\ast = \BreakEven$.

\section{The prefix probe}
\label{sec:probe}

Offline, the language model was given the first fraction $f$ of each
multi-step problem's reference transcript, for $f \in \{0.25, 0.5, 0.75, 0.9,
1\}$, and its draft was compared with the reply to the full transcript. The 80
seeded items were run twice: once sampling at temperature 0.5, once greedy.

\textbf{The gate.} The probe is meaningful only if the model can answer the
full problem. From the full transcript, its final answer is correct on
\SampledGate{} of items when sampling, which fails the registered 30\% gate.
Greedy, it is correct on \GreedyGate, which passes. The usability curve below
is therefore the greedy one; the live pipeline itself replies at temperature
0.5.

\textbf{The curve.} Greedy usability is \GreedyFifty{} at half the utterance,
\GreedySeventyFive{} at three-quarters (\GreedySeventyFiveK{} items) and
\GreedyNinety{} at 0.9. When sampling, usability at three-quarters is
\SampledSeventyFive.

\textbf{The registered test.} Before the run we registered that if greedy
usability at three-quarters of the utterance fell below the break-even, the
value side would be closed on this corpus. It is \GreedySeventyFive, against
$p^\ast = \BreakEven$: the test fires, and the value side is closed. At 0.9
the point estimate clears the break-even. That figure was exploratory, not
the registered test, and its interval reaches below the break-even.

\textbf{Why.} The problems ask their question last. A draft can anticipate the
answer only once the question has been heard, and by then the utterance is
nearly over.

\section{Where a draft still fits}
\label{sec:fstar}

\begin{figure}[t]
  \centering
  \includegraphics[width=0.85\textwidth]{paper_fig4_usable_vs_window.pdf}
  \caption{Draft usability rises with the fraction heard; the window to finish
  a draft falls. The last fraction at which a 10-token draft fits, as
  predicted and as measured.}
  \label{fig:th-window}
\end{figure}

Usability rises with the fraction heard, and the time left to finish a draft
falls (Figure~\ref{fig:th-window}). We registered a window model: the speech
left after the decode of what has been heard. From it we predicted the last
fraction at which a 10-token draft still fits, $f^\ast = 0.858$. Measured on
consistent clocks it is \Fstar{} ($n=\FstarN$ items in the interquartile range
of duration). The first estimate, on mixed clocks, was higher, and its
correction is on record.

At $f^\ast$, greedy usability on the probe's curve is \FstarUsable, within
sampling error of the break-even \BreakEven. This is the crux of the negative.
Where a draft fits, it is rarely worth issuing. Where it would be worth
issuing, it has stopped fitting. The two curves meet at the edge of
break-even; that usability there is above break-even is not shown.

\section{Could a trigger tell?}
\label{sec:triggers}

A trigger that predicted the remaining speech could time drafts better than a
fixed fraction. Two were evaluated against whether more than 900\,ms of speech
remained: a semantic trigger on the partial transcript, and an acoustic
endpoint-anticipation model. Over \TriggerPoints{} decision points from
\TriggerUtts{} utterances of the development set (and \TriggerUnindexed{}
points with no utterance index), their AUCs are \TsemAUC{} and \EpaAUC. Neither
is distinguishable from chance on this set, and the intervals are too wide to
call either useless. The question is open beyond ``not detectable''. Reading
the triggers on the multi-step set was registered not to happen once the
probe closed the value side, and it was not done. The acoustic model is also
disqualified on time alone: it runs slower than real time on this device
(\needsclaim{acoustic trigger, seconds per second of audio}{CORRECTION
(2026-09-21) to the EPA entry above}).

\section{Drafts in the live pipeline}
\label{sec:live-thinker}

Live drafts were no better. In the first live comparison, on the development
set and before the probe existed, drafts that restarted on every new partial
matched the final reply's first sentence in
none of the turns where one existed
(\needsclaim{first-sentence matches}{SPEC-ALWAYS-PG vs SPEC-CONTINUE: the
prediction scored}). Successive drafts shared few of their tokens
(\needsclaim{token survival between candidates}{same entry}). Restarting also
made TTFA worse than not speculating
(\needsclaim{SPEC-ALWAYS-PG minus REACTIVE}{same entry}).

In the last comparison, a continuing draft was interleaved with the reactive
pipeline on multi-step problems. Each draft produced \SpecTokens{} tokens per
turn, the full budget, and the per-item median share of its tokens accepted
was \SpecAccepted{} ($n=\SpecItems$ items). TTFA changed by \SpecTTFAr{}
($n=\SpecN$ items, positive is slower) at \SpecEnergyr{} per turn. This
comparison is descriptive. Its interleaved design broke at the first
voice-activity split, and its scorer was written after the run.

\section{Summary}

On this task, with this model, the thinker's speculation has a measured cost
(Chapter~\ref{ch:cost}) and no measured value. The negative is specific. A
draft becomes usable only near the end of the utterance, after the question
has been heard. By then the window to finish it has closed or nearly closed,
and neither trigger we evaluated could be shown to predict the end. The
result may differ on tasks whose questions come early. The allocator in
Chapter~\ref{ch:allocator} reads the measured usability curve, so on such a
task its rule would admit drafts wherever the curve clears break-even within
the window.
